\section{附录 H：符号 \texorpdfstring{$\zeta$}{zeta} 生成函数与算术结构接口（扩展，非主链）}
\label{app:zeta-interface}

本附录给出一个严格但有限的接口：说明“素数因子结构如何在周期轨道计数/$\zeta$ 生成函数的反演公式中出现”，并不声称给出素数分布定理或素数生成算法。

\subsection{符号动力系统的 Artin--Mazur \texorpdfstring{$\zeta$}{zeta} 函数}

对有限状态移位（SFT），周期点计数 $\card{\mathrm{Fix}(T^n)}$ 定义 $\zeta$ 生成函数
\[
\zeta_T(z):=\exp\Bigl(\sum_{n\ge 1}\frac{\card{\mathrm{Fix}(T^n)}}{n}z^n\Bigr).
\]
对邻接矩阵 $A$，有封闭形式
\[
\zeta_T(z)=\frac{1}{\det(I-zA)}.
\]
对黄金均值移位（禁词 $11$），邻接矩阵可取
\[
A=\begin{pmatrix}1&1\\1&0\end{pmatrix},
\quad\Rightarrow\quad
\det(I-zA)=1-z-z^2,
\]
因此
\[
\zeta_T(z)=\frac{1}{1-z-z^2}.
\]

\subsection{原始周期轨道与 Möbius 反演}

设 $P(n)$ 为长度为 $n$ 的原始周期轨道数，则经典关系为
\[
\card{\mathrm{Fix}(T^n)}=\sum_{d\mid n} d\,P(d).
\]
由 Möbius 反演得
\[
P(n)=\frac{1}{n}\sum_{d\mid n}\mu(d)\,\card{\mathrm{Fix}(T^{n/d})},
\]
其中 $\mu$ 为 Möbius 函数。该式表明：素因子结构通过约数格与 Möbius 函数进入“原始—非原始”分解。这提供一个可检验的算术接口：一旦周期点计数已知，原始轨道计数由包含素因子信息的反演给出。

\subsection{与主文协议的关系（解释层）}

主文的“周期近似”在参数逼近层面产生有限周期结构；若进一步将这些周期结构编码为 SFT 或其有限近似，则 $\zeta$ 生成函数与 Möbius 反演为“从周期计数提取原始结构”的标准工具。将该原始结构解释为“素数生成”需额外语义映射与外部算术对象选择，属于解释层，不进入闭层定理链。

